\centerline{\bf \Huge Теорема Без у...}
\centerline{\bf \Huge }

\begin{flushright}
        \scriptsize{...без ума}
\end{flushright}

\textbf{Определение.} Многочлен $A$ делится на ненулевой многочлен $B$, если существует многочлен $Q$, называемый частным такой, что $A=B \cdot Q$.

\textbf{Определение.} Разделить многочлен $A$ на ненулевой многочлен $B$ с остатком это найти многочлены $Q, R$ такие, что выполнено равенство $A=B \cdot Q+R$, причем $\operatorname{deg} R<\operatorname{deg} B$ или $R=0$. Многочлен $Q$ называется неполным частным, многочлен $R$ называется остатком.

\textbf{Теорема Безу}. Для любого многочлена $P(x)$ и любого действительного числа $a$ верно:

$$
P(x) \equiv P(a) \ (mod \ (x-a))
$$

или же 

$$
P(x)=H(x) \cdot(x-a)+P(a) .
$$


\textbf{Важное следствие.} Число $a$ является корнем многочлена $P(x)$ тогда и только тогда, когда $P(x)$ делится на $(x-a)$.

\textbf{Теорема о рациональном корне.} Если многочлен с целыми коэффициентами

$$
a_n x^n+a_{n-1}x_{n-1}+\ldots+a_{1} x+a_0
$$


имеет рациональный корень $p / q$ (дробь несократима), то 

\begin{enumerate}
    \item $a_n \ \vdots \ q$
    \item $a_0 \ \vdots \ p$
\end{enumerate}


\newpage

\hspace{3mm}

\centerline{\bf \Huge Теорема Без у}

\begin{flushright}
        \scriptsize{какая-то из этих задач будет в варианте Q.}
\end{flushright}

\problem{1}
Приведенный кубический многочлен имеет три положительных корня, каждый из которых не превосходит 2. Докажите, что сумма коэффициентов этого многочлена по модулю не превосходит 1.

\problem{2}
Найдите остаток при делении многочлена $x^{100}-8 x^{97}-5 x^{17}+10 x^{16}+x^2-2 x+1$ на $x^2-3 x+2$.

\problem{3}
Числа $x, y, z$ удовлетворяют равенству $x+y+z-2(x y+y z+x z)+4 x y z= \dfrac{1}{2}$. Докажите, что хотя бы одно из них равно $\dfrac{1}{2}$.

\problem{4}
Олег попросил Эльзяту назвать любой многочлен $P(x)$ с натуральными коэффициентами, а затем назвал такое число $k$, что $P(k), P(k+1), \ldots, P(k+2024)$ являются составными числами. Эльзята подумала, что её брату просто повезло и придумала ещё один многочлен. Но Олег и тут смог найти подходящее число $k$. Правда ли, что Олег всегда сможет находить такое число, какой бы многочлен не загадывала Эльзята или ему просто дважды повезло?

\problem{5}
Разложите на множители выражение $(a + b + c)^3 - a^3 - b^3 - c^3$


\problem{6}
Многочлен $P(x)=x^3+a x^2+b x+c$ имеет три различных действительных корня, а многочлен $P(Q(x))$, где $Q(x)=x^2+x+2025$, действительных корней не имеет. Докажите, что $P(2025)>\frac{1}{64}$.

\problem{7}
0) \textit{Вспомогательная лемма к вспомогательной лемме.} Докажите, что у любого многочлена степени $n$ не более чем $n$ корней.

\noindent
1) \textit{Вспомогательная лемма.} Многочлен степени $n$ не может принимать одно и то же значение в более чем $n$ различных точках.

\noindent
2) Олежа выписал на доске 1000 различных чисел. Затем он увеличил все числа на один. Оказалось, что произведение чисел не изменилось. Затем он повторил операцию, и произведение чисел снова не изменилось. Какое наибольшее количество раз он мог повторить эту операцию, чтобы произведение чисел оставалось постоянным?